\section{Preschemes and morphisms of preschemes}
\label{section:I.2}

\subsection{Definition of preschemes.}
\label{subsection:I.2.1}

\begin{env}[2.1.1]
\label{I.2.1.1}
Given a ringed space $(X,\mathscr{O}_X)$, we say that an open subset
$V$ of $X$ is an \emph{affine open set} if the ringed space
$(V,\mathscr{O}_X|V)$ is an affine scheme
(\hyperref[I.1.7.1]{1.7.1}).
\end{env}

\begin{definition}[2.1.2]
\label{I.2.1.2}
We call \emph{prescheme} a ringed space $(X,\mathscr{O}_X)$ such that every
point of $X$ has an affine open neighbourhood.
\end{definition}

\oldpage[I]{98}

\begin{proposition}[2.1.3]
\label{I.2.1.3}
If $(X,\mathscr{O}_X)$ is a prescheme, the affine open sets form a base for
the topology of $X$.
\end{proposition}

Indeed, if $V$ is an arbitrary open neighbourhood of $x\in X$, there exists
by hypothesis an open neighbourhood $W$ of $x$ such that
$(W,\mathscr{O}_X|W)$ is an affine scheme; let $A$ denote its ring.
In the space $W$, $V\cap W$ is an open neighbourhood of $x$; hence there
exists $f\in A$ such that $D(f)$ is an open neighbourhood of $x$ contained in
$V\cap W$ (\hyperref[I.1.1.10]{1.1.10 (i)}). The ringed space
$(D(f),\mathscr{O}_X|D(f))$ is then an affine scheme with ring isomorphic to
$A_f$ (\hyperref[I.1.3.6]{1.3.6}), whence the proposition.

\begin{proposition}[2.1.4]
\label{I.2.1.4}
The underlying space of a prescheme is a Kolmogorov space.
\end{proposition}

Indeed, if $x,y$ are two distinct points of a prescheme $X$, it is
obvious that there exists an open neighbourhood of one of these points not
containing the other if $x,y$ are not in one and the same affine open set;
and if they are in one and the same affine open set, this follows from
(\hyperref[I.1.1.8]{1.1.8}).

\begin{proposition}[2.1.5]
\label{I.2.1.5}
If $(X,\mathscr{O}_X)$ is a prescheme, every irreducible closed subset of
$X$ has one and only one generic point, and the map
$x\longmapsto\overline{\{x\}}$ is therefore a bijection of $X$ onto the set
of its irreducible closed subsets.
\end{proposition}

Indeed, if $Y$ is an irreducible closed subset of $X$ and $y\in Y$, and if
$U$ is an affine open neighbourhood of $y$ in $X$, $U\cap Y$ is dense in $Y$
and is irreducible
(\hyperref[0.2.1.1]{0, 2.1.1} and
\hyperref[0.2.1.4]{2.1.4}); hence
(\hyperref[I.1.1.14]{1.1.14}) $U\cap Y$ is the closure in $U$ of a
point $x$, and consequently $Y\subset\overline{U}$ is the closure of $x$ in
$X$. The uniqueness of the generic point of $Y$ follows from
(\hyperref[I.2.1.4]{2.1.4}) and from
(\hyperref[0.2.1.3]{0, 2.1.3}).

\begin{env}[2.1.6]
\label{I.2.1.6}
If $Y$ is an irreducible closed subset of $X$ and $y$ its generic point,
the local ring $\mathscr{O}_y$ is also denoted by $\mathscr{O}_{X/Y}$ and is
called the \emph{local ring of $X$ along $Y$}, or the \emph{local ring of $Y$
in $X$}.

If $X$ itself is irreducible and $x$ is its generic point, we also say
that $\mathscr{O}_x$ is the \emph{ring of rational functions on
$X$} (cf. \hyperref[section:I.7]{§~7}).
\end{env}

\begin{proposition}[2.1.7]
\label{I.2.1.7}
If $(X,\mathscr{O}_X)$ is a prescheme, then for every open subset $U$ of
$X$, the ringed space $(U,\mathscr{O}_X|U)$ is a prescheme.
\end{proposition}

This follows at once from Definition (\hyperref[I.2.1.2]{2.1.2}) and
Proposition (\hyperref[I.2.1.3]{2.1.3}).

We say that $(U,\mathscr{O}_X|U)$ is the prescheme \emph{induced} on $U$ by
$(X,\mathscr{O}_X)$, or the \emph{restriction} of $(X,\mathscr{O}_X)$ to $U$.

\begin{env}[2.1.8]
\label{I.2.1.8}
A prescheme $(X,\mathscr{O}_X)$ is said to be \emph{irreducible} (resp.
\emph{connected}) if the underlying space $X$ is irreducible (resp.
connected). A prescheme is said to be \emph{integral} if it is
\emph{irreducible and reduced} (cf. (\hyperref[I.5.1.4]{5.1.4})). A
prescheme $(X,\mathscr{O}_X)$ is said to be \emph{locally integral} if every
$x\in X$ has an open neighbourhood $U$ such that the prescheme induced on $U$
by $(X,\mathscr{O}_X)$ is integral.
\end{env}

\subsection{Morphisms of preschemes.}
\label{subsection:I.2.2}
\label{I.2.2}

\begin{definition}[2.2.1]
\label{I.2.2.1}
Given two preschemes $(X,\mathscr{O}_X)$, $(Y,\mathscr{O}_Y)$, we call
\emph{morphism (of preschemes)} of $(X,\mathscr{O}_X)$ into
$(Y,\mathscr{O}_Y)$ every morphism of ringed spaces $(\psi,\theta)$ such
that, for every $x\in X$, $\theta_x^\sharp$ is a local homomorphism:
$\mathscr{O}_{\psi(x)}\longrightarrow\mathscr{O}_x$.
\end{definition}

\oldpage[I]{99}

By passing to the quotients, the map
$\theta_x^\sharp:\mathscr{O}_{\psi(x)}\to\mathscr{O}_x$ therefore gives a
monomorphism $\theta^x:k(\psi(x))\to k(x)$, which makes it possible to
regard $k(x)$ as an \emph{extension} of the field $k(\psi(x))$.

\begin{env}[2.2.2]
\label{I.2.2.2}
The composite $(\psi'',\theta'')$ of two morphisms of preschemes
$(\psi,\theta)$ and $(\psi',\theta')$ is again a morphism of preschemes,
as follows from the formula
${\theta''}^\sharp=\theta^\sharp\circ\psi^*({\theta'}^\sharp)$
(\hyperref[0.3.5.5]{0, 3.5.5}). We conclude that the preschemes form
a \emph{category}; in accordance with the general notation, we shall write
$\operatorname{Hom}(X,Y)$ for the set of morphisms of a prescheme $X$ into
a prescheme $Y$.
\end{env}

\begin{example}[2.2.3]
\label{I.2.2.3}
If $U$ is an open subset of $X$, the canonical injection
(\hyperref[0.4.1.2]{0, 4.1.2}) of the induced prescheme
$(U,\mathscr{O}_X|U)$ into $(X,\mathscr{O}_X)$ is a morphism of
preschemes; it is moreover a \emph{monomorphism} of ringed spaces (and
\emph{a fortiori} a monomorphism of preschemes), as follows at once
from (\hyperref[0.4.1.1]{0, 4.1.1}).
\end{example}

\begin{proposition}[2.2.4]
\label{I.2.2.4}
Let $(X,\mathscr{O}_X)$ be a prescheme and $(S,\mathscr{O}_S)$ an affine
scheme with ring $A$. There is a canonical one-to-one correspondence between
the morphisms of the prescheme $(X,\mathscr{O}_X)$ into the prescheme
$(S,\mathscr{O}_S)$ and the homomorphisms of $A$ into the ring
$\Gamma(X,\mathscr{O}_X)$.
\end{proposition}

Let us first note that, if $(X,\mathscr{O}_X)$ and $(Y,\mathscr{O}_Y)$ are
two arbitrary ringed spaces, a morphism $(\psi,\theta)$ of
$(X,\mathscr{O}_X)$ into $(Y,\mathscr{O}_Y)$ canonically defines a ring
homomorphism
$\Gamma(\theta):\Gamma(Y,\mathscr{O}_Y)\to
\Gamma(Y,\psi_*(\mathscr{O}_X))=\Gamma(X,\mathscr{O}_X)$. In the case under
consideration, everything comes down to showing that an arbitrary
homomorphism $\varphi:A\to\Gamma(X,\mathscr{O}_X)$ is of the form
$\Gamma(\theta)$ for one and only one $\theta$. Now, by hypothesis there is
a covering $(V_\alpha)$ of $X$ by affine open sets; by composing $\varphi$
with the restriction homomorphism
$\Gamma(X,\mathscr{O}_X)\to
\Gamma(V_\alpha,\mathscr{O}_X|V_\alpha)$, we obtain a homomorphism
$\varphi_\alpha:A\to\Gamma(V_\alpha,\mathscr{O}_X|V_\alpha)$, which
corresponds to a unique morphism $(\psi_\alpha,\theta_\alpha)$ of the
prescheme $(V_\alpha,\mathscr{O}_X|V_\alpha)$ into $(S,\mathscr{O}_S)$, by
virtue of (\hyperref[I.1.7.3]{1.7.3}). Moreover, for every pair of
indices $(\alpha,\beta)$, every point of $V_\alpha\cap V_\beta$ has an
affine open neighbourhood $W$ contained in $V_\alpha\cap V_\beta$
(\hyperref[I.2.1.3]{2.1.3}); it is clear that by composing
$\varphi_\alpha$ and $\varphi_\beta$ with the restriction homomorphisms to
$W$, we obtain the same homomorphism
$\Gamma(S,\mathscr{O}_S)\to\Gamma(W,\mathscr{O}_X|W)$; hence, by virtue of
the relations $(\theta_\alpha^\sharp)_x=(\varphi_\alpha)_x$ for every
$x\in V_\alpha$ and every $\alpha$ (\hyperref[I.1.6.1]{1.6.1}), the
restrictions to $W$ of the morphisms $(\psi_\alpha,\theta_\alpha)$ and
$(\psi_\beta,\theta_\beta)$ coincide. We conclude that there is one and only
one morphism of ringed spaces
$(\psi,\theta):(X,\mathscr{O}_X)\to(S,\mathscr{O}_S)$ whose
restriction to each $V_\alpha$ is $(\psi_\alpha,\theta_\alpha)$, and it is
clear that this morphism is a morphism of preschemes and is such that
$\Gamma(\theta)=\varphi$.

Let $u:A\to\Gamma(X,\mathscr{O}_X)$ be a ring homomorphism, and let
$v=(\psi,\theta)$ be the corresponding morphism
$(X,\mathscr{O}_X)\to(S,\mathscr{O}_S)$. For every $f\in A$, we have
\begin{equation}
  \psi^{-1}(D(f))=X_{u(f)}
  \tag{2.2.4.1}\label{I.2.2.4.1}
\end{equation}
with the notation of (\hyperref[0.5.5.2]{0, 5.5.2}) relative to the
locally free sheaf $\mathscr{O}_X$. Indeed, it suffices to verify this
formula when $X$ itself is affine, and then it is none other than
(\hyperref[I.1.2.2.2]{1.2.2.2}).

\begin{proposition}[2.2.5]
\label{I.2.2.5}
Under the hypotheses of (\hyperref[I.2.2.4]{2.2.4}), let
$\varphi:A\to\Gamma(X,\mathscr{O}_X)$ be a ring homomorphism,
$f:(X,\mathscr{O}_X)\to(S,\mathscr{O}_S)$ the corresponding morphism of
preschemes, $\mathscr{G}$ (resp. $\mathscr{F}$) a quasi-coherent
$\mathscr{O}_X$-module (resp. $\mathscr{O}_S$-module), and
let $M=\Gamma(S,\mathscr{F})$. There is then a
\oldpage[I]{100}
canonical one-to-one correspondence between the $f$-morphisms
$\mathscr{F}\to\mathscr{G}$ (\hyperref[0.4.4.1]{0, 4.4.1}) and the
$A$-homomorphisms
$M\to(\Gamma(X,\mathscr{G}))_{[\varphi]}$.
\end{proposition}

Indeed, by reasoning as in (\hyperref[I.2.2.4]{2.2.4}), we are at once
reduced to the case where $X$ is affine, and the proposition then follows
from (\hyperref[I.1.6.3]{1.6.3}) and (\hyperref[I.1.3.8]{1.3.8}).

\begin{env}[2.2.6]
\label{I.2.2.6}
A morphism of preschemes
$(\psi,\theta):(X,\mathscr{O}_X)\to(Y,\mathscr{O}_Y)$ is said to be
\emph{open} (resp. \emph{closed}) if for every open subset $U$ of $X$ (resp.
every closed subset $F$ of $X$), $\psi(U)$ is open in $Y$ (resp. $\psi(F)$
is closed in $Y$). We say that $(\psi,\theta)$ is \emph{dominant} if
$\psi(X)$ is dense in $Y$, and \emph{surjective} if $\psi$ is surjective.
Note that these conditions involve only the continuous map $\psi$.
\end{env}

\begin{proposition}[2.2.7]
\label{I.2.2.7}
Let
\[
  f=(\psi,\theta):(X,\mathscr{O}_X)\to(Y,\mathscr{O}_Y),
  \qquad
  g=(\psi',\theta'):(Y,\mathscr{O}_Y)\to(Z,\mathscr{O}_Z)
\]
be two morphisms of preschemes.
\begin{enumerate}
  \item[\rm(i)] If $f$ and $g$ are both open (resp. closed, dominant,
    surjective), so is $g\circ f$.
  \item[\rm(ii)] If $f$ is surjective and $g\circ f$ is closed, $g$ is
    closed.
  \item[\rm(iii)] If $g\circ f$ is surjective, $g$ is surjective.
\end{enumerate}
\end{proposition}

Assertions (i) and (iii) are obvious. Set
$g\circ f=(\psi'',\theta'')$. If $F$ is closed in $Y$,
$\psi^{-1}(F)$ is closed in $X$, hence
$\psi''(\psi^{-1}(F))$ is closed in $Z$; but since $\psi$ is surjective,
$\psi(\psi^{-1}(F))=F$, hence
$\psi''(\psi^{-1}(F))=\psi'(F)$, which proves (ii).

\begin{proposition}[2.2.8]
\label{I.2.2.8}
Let $f=(\psi,\theta)$ be a morphism
$(X,\mathscr{O}_X)\to(Y,\mathscr{O}_Y)$ and $(U_\alpha)$ an open covering
of $Y$. For $f$ to be open (resp. closed, surjective, dominant), it is
necessary and sufficient that its restriction to each of the induced
preschemes $(\psi^{-1}(U_\alpha),\mathscr{O}_X|\psi^{-1}(U_\alpha))$,
regarded as a morphism of this induced prescheme into the induced prescheme
$(U_\alpha,\mathscr{O}_Y|U_\alpha)$, be open (resp. closed, surjective,
dominant).
\end{proposition}

The proposition follows at once from the definitions, taking into account
the fact that a subset $F$ of $Y$ is closed (resp. open, dense) in $Y$ if
and only if each of the sets $F\cap U_\alpha$ is closed (resp. open, dense)
in $U_\alpha$.

\begin{env}[2.2.9]
\label{I.2.2.9}
Let $(X,\mathscr{O}_X)$, $(Y,\mathscr{O}_Y)$ be two preschemes; suppose
that $X$ and $Y$ have the same \emph{finite} number of irreducible
components $X_i$ (resp. $Y_i$) $(1\leq i\leq n)$; let $\xi_i$ (resp.
$\eta_i$) be the generic point of $X_i$ (resp. $Y_i$)
(\hyperref[I.2.1.5]{2.1.5}). A morphism
\[
  f=(\psi,\theta):(X,\mathscr{O}_X)\to(Y,\mathscr{O}_Y)
\]
is said to be \emph{birational} if, for every $i$,
$\psi^{-1}(\eta_i)=\{\xi_i\}$ and
$\theta_{\xi_i}^{\sharp}:\mathscr{O}_{\eta_i}\to\mathscr{O}_{\xi_i}$ is an
\emph{isomorphism}. It is clear that a birational morphism is
\emph{dominant} (\hyperref[0.2.1.8]{0, 2.1.8}), and hence surjective if
it is also closed.
\end{env}

\begin{notation}[2.2.10]
\label{I.2.2.10}
Throughout the rest of this work, and when this is not likely to cause
confusion, we shall \emph{suppress} from the notation for a prescheme (resp.
a morphism) the structure sheaf (resp. the morphism of structure sheaves).
If $U$ is an open subset of the underlying space of a prescheme $X$, when we
speak of $U$ as a prescheme, it will always be the prescheme induced on $U$.
\end{notation}

\subsection{Gluing of preschemes}
\label{subsection:I.2.3}

\begin{env}[2.3.1]
\label{I.2.3.1}
\oldpage[I]{101}
It follows from Definition (\hyperref[I.2.1.2]{2.1.2}) that every ringed
space obtained by \emph{gluing} preschemes
(\hyperref[0.4.1.6]{0, 4.1.6}) is again a prescheme. In particular,
since every prescheme has by definition a covering consisting of affine
open sets, we see that every prescheme can be obtained by
\emph{gluing affine schemes}.
\end{env}

\begin{example}[2.3.2]
\label{I.2.3.2}
Let $K$ be a field, $B=K[s]$, $C=K[t]$ two polynomial rings in one
indeterminate over $K$, and let $X_1=\operatorname{Spec}(B)$,
$X_2=\operatorname{Spec}(C)$, which are two isomorphic affine schemes. In
$X_1$ (resp. $X_2$), let $U_{12}$ (resp. $U_{21}$) be the affine open set
$D(s)$ (resp. $D(t)$), whose ring $B_s$ (resp. $C_t$) consists of the
rational fractions of the form $f(s)/s^m$ (resp. $g(t)/t^n$) with $f\in B$
(resp. $g\in C$). Let $u_{12}$ be the isomorphism of preschemes
$U_{21}\to U_{12}$ corresponding (\hyperref[I.2.2.4]{2.2.4}) to
the isomorphism of $B_s$ onto $C_t$ which maps $f(s)/s^m$ to the rational
fraction $f(1/t)/(1/t^m)$. We can glue $X_1$ and $X_2$ along
$U_{12}$ and $U_{21}$ by means of $u_{12}$, since there is obviously no
gluing condition. We shall meet the prescheme $X$ thus obtained again later,
as a special case of a general method of construction
(\hyperref[II.2.4.3]{II, 2.4.3}). Let us only show here that $X$
\emph{is not an affine scheme}; this will follow from the fact that the ring
$\Gamma(X,\mathscr{O}_X)$ is \emph{isomorphic} to $K$, and therefore has a
spectrum consisting of a single point. Indeed, a section of $\mathscr{O}_X$
over $X$ has a restriction over $X_1$ (resp. $X_2$), identified with an
affine open set of $X$, which is a polynomial $f(s)$ (resp. $g(t)$), and it
follows from the definitions that we must have $g(t)=f(1/t)$, which is
possible only if $f=g\in K$.
\end{example}

\subsection{Local schemes}
\label{subsection:I.2.4}

\begin{env}[2.4.1]
\label{I.2.4.1}
We call \emph{local scheme} an affine scheme whose ring $A$ is local;
there is then in $X=\operatorname{Spec}(A)$ a single \emph{closed point}
$a$, and for every other point $b\in X$, we have
$a\in\overline{\{b\}}$ (\hyperref[I.1.1.7]{1.1.7}).
\end{env}

For every prescheme $Y$ and every point $y\in Y$, the local scheme
$\operatorname{Spec}(\mathscr{O}_y)$ is called the \emph{local scheme of $Y$
at the point $y$}. Let $V$ be an affine open set of $Y$ containing $y$, and
$B$ the ring of the affine scheme $V$; $\mathscr{O}_y$ is canonically
identified with $B_y$ (\hyperref[I.1.3.4]{1.3.4}), and the canonical
homomorphism $B\to B_y$ consequently corresponds
(\hyperref[I.1.6.1]{1.6.1}) to a morphism of
preschemes $\operatorname{Spec}(\mathscr{O}_y)\to V$. If we compose this
morphism with the canonical injection $V\to Y$, we therefore obtain a
morphism $\operatorname{Spec}(\mathscr{O}_y)\to Y$, which is
\emph{independent} of the affine open set $V$ (containing $y$) chosen:
indeed, if $V'$ is a second affine open set containing $y$, there exists a
third affine open set $W$ containing $y$ and such that $W\subset V\cap V'$
(\hyperref[I.2.1.3]{2.1.3}); we may therefore restrict ourselves to the
case where $V\subset V'$, and if $B'$ is the ring of $V'$, everything comes
down to observing that the diagram
\phantomsection\label{I.2.4.1.diagram}
\[
  \xymatrix{
    B'\ar[rr]\ar[dr] & & B\ar[dl]\\
    & \mathscr{O}_y
  }
\]
is commutative (\hyperref[0.1.5.1]{0, 1.5.1}). The morphism
\[
  \operatorname{Spec}(\mathscr{O}_y)\to Y
\]
thus defined is called \emph{canonical}.

\begin{proposition}[2.4.2]
\label{I.2.4.2}
\oldpage[I]{102}
Let $(Y,\mathscr{O}_Y)$ be a prescheme; for every $y\in Y$, let
$(\psi,\theta)$ be the canonical morphism
$(\operatorname{Spec}(\mathscr{O}_y),\widetilde{\mathscr{O}}_y)
\to(Y,\mathscr{O}_Y)$. Then $\psi$ is a homeomorphism of
$\operatorname{Spec}(\mathscr{O}_y)$ onto the subspace $S_y$ of $Y$
consisting of the $z$ such that $y\in\overline{\{z\}}$ (in other words, of
the \emph{generizations} of $y$ (\hyperref[0.2.1.2]{0, 2.1.2}));
moreover, if $z=\psi(\mathfrak{p})$,
$\theta_z^{\sharp}:\mathscr{O}_z\to(\mathscr{O}_y)_{\mathfrak{p}}$ is an
isomorphism; $(\psi,\theta)$ is therefore a monomorphism of ringed spaces.
\end{proposition}

Since the unique closed point $a$ of $\operatorname{Spec}(\mathscr{O}_y)$
lies in the closure of every point of this space, and since $\psi(a)=y$, the
image of $\operatorname{Spec}(\mathscr{O}_y)$ under the continuous map $\psi$
is contained in $S_y$. Since $S_y$ is contained in every affine open set
containing $y$, we may reduce to the case where $Y$ is an affine scheme; but
in that case the proposition follows from
(\hyperref[I.1.6.2]{1.6.2}).

\emph{We therefore see (\hyperref[I.2.1.5]{2.1.5}) that there is a
one-to-one correspondence between $\operatorname{Spec}(\mathscr{O}_y)$ and
the set of irreducible closed subsets of $Y$ containing $y$.}

\begin{corollary}[2.4.3]
\label{I.2.4.3}
For $y\in Y$ to be the generic point of an irreducible component of
$Y$, it is necessary and sufficient that the only prime ideal of the local
ring $\mathscr{O}_y$ be its maximal ideal (in other words, that
$\mathscr{O}_y$ be of dimension zero).
\end{corollary}

\begin{proposition}[2.4.4]
\label{I.2.4.4}
Let $(X,\mathscr{O}_X)$ be a local scheme with ring $A$, $a$ its unique
closed point, and $(Y,\mathscr{O}_Y)$ a prescheme. Every morphism
$u=(\psi,\theta):(X,\mathscr{O}_X)\to(Y,\mathscr{O}_Y)$ factors uniquely
as $X\to\operatorname{Spec}(\mathscr{O}_{\psi(a)})\to Y$, where the second
arrow denotes the canonical morphism and the first corresponds to a local
homomorphism $\mathscr{O}_{\psi(a)}\to A$. This establishes a canonical
one-to-one correspondence between the set of morphisms
$(X,\mathscr{O}_X)\to(Y,\mathscr{O}_Y)$ and the set of local homomorphisms
$\mathscr{O}_y\to A$ $(y\in Y)$.
\end{proposition}

Indeed, for every $x\in X$, we have $a\in\overline{\{x\}}$, hence
$\psi(a)\in\overline{\{\psi(x)\}}$, which proves that $\psi(X)$ is contained
in every affine open set containing $\psi(a)$. We may therefore reduce to
the case where $(Y,\mathscr{O}_Y)$ is an affine scheme with ring $B$, and we
then have $u=({}^a\varphi,\widetilde{\varphi})$, where
$\varphi\in\operatorname{Hom}(B,A)$ (\hyperref[I.1.7.3]{1.7.3}).
Moreover, we have $\varphi^{-1}(\mathfrak{j}_a)=\mathfrak{j}_{\psi(a)}$, and
consequently the image under $\varphi$ of every element of
$B-\mathfrak{j}_{\psi(a)}$ is invertible in the local ring $A$; the
factorization in the statement therefore follows from the universal property
of rings of fractions (\hyperref[0.1.2.4]{0, 1.2.4}). Conversely, to
every local homomorphism $\mathscr{O}_y\to A$ there corresponds a unique
morphism $(\psi,\theta):X\to\operatorname{Spec}(\mathscr{O}_y)$ such that
$\psi(a)=y$ (\hyperref[I.1.7.3]{1.7.3}), and by composing it with the
canonical morphism $\operatorname{Spec}(\mathscr{O}_y)\to Y$, we obtain a
morphism $X\to Y$, which completes the proof of the proposition.

\begin{env}[2.4.5]
\label{I.2.4.5}
The affine schemes whose ring is a \emph{field} $K$ have an underlying
space consisting of a single point. If $A$ is a local ring with maximal
ideal $\mathfrak{m}$, every local homomorphism $A\to K$ has kernel equal to
$\mathfrak{m}$, and therefore factors as
$A\to A/\mathfrak{m}\to K$, where the second arrow is a monomorphism. The
morphisms $\operatorname{Spec}(K)\to\operatorname{Spec}(A)$ therefore
correspond one-to-one to the monomorphisms of fields
$A/\mathfrak{m}\to K$.
\end{env}

Let $(Y,\mathscr{O}_Y)$ be a prescheme; for every $y\in Y$ and every ideal
$\mathfrak{a}_y$ of $\mathscr{O}_y$, the canonical homomorphism
$\mathscr{O}_y\to\mathscr{O}_y/\mathfrak{a}_y$ defines a morphism
$\operatorname{Spec}(\mathscr{O}_y/\mathfrak{a}_y)
\to\operatorname{Spec}(\mathscr{O}_y)$; if we compose it with the canonical
morphism $\operatorname{Spec}(\mathscr{O}_y)\to Y$, we obtain a morphism
$\operatorname{Spec}(\mathscr{O}_y/\mathfrak{a}_y)\to Y$, also called
\emph{canonical}. For $\mathfrak{a}_y=\mathfrak{m}_y$, that is to say
$\mathscr{O}_y/\mathfrak{a}_y=k(y)$, Proposition
(\hyperref[I.2.4.4]{2.4.4}) therefore implies:

\begin{corollary}[2.4.6]
\label{I.2.4.6}
\oldpage[I]{103}
Let $(X,\mathscr{O}_X)$ be a local scheme whose ring is a field $K$,
$\xi$ the unique point of $X$, and $(Y,\mathscr{O}_Y)$ a prescheme. Every
morphism $u:(X,\mathscr{O}_X)\to(Y,\mathscr{O}_Y)$ factors uniquely as
$X\to\operatorname{Spec}(k(\psi(\xi)))\to Y$, where the second arrow denotes
the canonical morphism and the first corresponds to a monomorphism
$k(\psi(\xi))\to K$. This establishes a canonical one-to-one correspondence
between the set of morphisms $(X,\mathscr{O}_X)\to(Y,\mathscr{O}_Y)$ and
the set of monomorphisms $k(y)\to K$ $(y\in Y)$.
\end{corollary}

\begin{corollary}[2.4.7]
\label{I.2.4.7}
For every $y\in Y$, every canonical morphism
$\operatorname{Spec}(\mathscr{O}_y/\mathfrak{a}_y)\to Y$ is a monomorphism
of ringed spaces.
\end{corollary}

We have already seen this when $\mathfrak{a}_y=0$
(\hyperref[I.2.4.2]{2.4.2}), and it suffices to apply
(\hyperref[I.1.7.5]{1.7.5}).

\begin{remark}[2.4.8]
\label{I.2.4.8}
Let $X$ be a local scheme and $a$ its unique closed point. Since every
affine open set containing $a$ is necessarily the whole of $X$, every
\emph{invertible} $\mathscr{O}_X$-module
(\hyperref[0.5.4.1]{0, 5.4.1}) is necessarily \emph{isomorphic to
$\mathscr{O}_X$} (or, as one also says, is \emph{trivial}). This property
does not hold in general for an arbitrary affine scheme
$\operatorname{Spec}(A)$; we shall see in Chapter~V that if $A$ is a normal
ring, it holds when $A$ is \emph{factorial}.
\end{remark}

\subsection{Preschemes over a prescheme}
\label{subsection:I.2.5}

\begin{definition}[2.5.1]
\label{I.2.5.1}
Given a prescheme $S$, we say that the data of a prescheme $X$ and of a
morphism of preschemes $\varphi:X\to S$ define a prescheme $X$ \emph{over
the prescheme $S$}, or an \emph{$S$-prescheme}; we say that $S$ is the
\emph{base prescheme} of the $S$-prescheme $X$. The morphism $\varphi$ is
called the \emph{structural morphism} of the $S$-prescheme $X$. When $S$ is
an affine scheme with ring $A$, we also say that $X$ endowed with $\varphi$
is a prescheme over the ring $A$ (or an $A$-prescheme).
\end{definition}

It follows from (\hyperref[I.2.2.4]{2.2.4}) that giving a prescheme
over a ring $A$ is equivalent to giving a prescheme
$(X,\mathscr{O}_X)$ whose structure sheaf $\mathscr{O}_X$ is a sheaf
of $A$-\emph{algebras}. \emph{An arbitrary prescheme can therefore always be
regarded in a unique way as a $\mathbf{Z}$-prescheme.}

If $\varphi:X\to S$ is the structural morphism of an $S$-prescheme $X$, we
say that a point $x\in X$ \emph{lies over a point $s\in S$} if
$\varphi(x)=s$. We say that $X$ \emph{dominates} $S$ if $\varphi$ is a
dominant morphism (\hyperref[I.2.2.6]{2.2.6}).

\begin{env}[2.5.2]
\label{I.2.5.2}
Let $X$, $Y$ be two $S$-preschemes; we say that a morphism of preschemes
$u:X\to Y$ is a \emph{morphism of preschemes over $S$} (or an
$S$-\emph{morphism}) if the diagram
\phantomsection\label{I.2.5.2.diagram}
\[
  \xymatrix{
    X\ar[rr]^u\ar[dr] & & Y\ar[dl]\\
    & S &
  }
\]
(where the oblique arrows are the structural morphisms) is commutative:
this implies that for every $s\in S$ and every $x\in X$ lying over $s$,
$u(x)$ must also lie over $s$.
\end{env}

This definition shows at once that the composite of two $S$-morphisms is an
$S$-morphism; the $S$-preschemes therefore form a \emph{category}.

We shall denote by $\operatorname{Hom}_S(X,Y)$ the set of $S$-morphisms of
an $S$-prescheme $X$ into an $S$-prescheme $Y$; the identity morphism of an
$S$-prescheme $X$ will be denoted by $1_X$.

When $S$ is an affine scheme with ring $A$, we shall also say
$A$-\emph{morphism} instead of $S$-morphism.

\begin{env}[2.5.3]
\label{I.2.5.3}
\oldpage[I]{104}
If $X$ is an $S$-prescheme and $v:X'\to X$ a morphism of preschemes, the
composite morphism $X'\xrightarrow{v}X\to S$ defines $X'$ as an
$S$-prescheme; in particular, every prescheme induced on an open set $U$ of
$X$ can be regarded as an $S$-prescheme by means of the canonical
injection.

If $u:X\to Y$ is an $S$-morphism of $S$-preschemes, the restriction of $u$
to every prescheme induced on an open set $U$ of $X$ is therefore an
$S$-morphism $U\to Y$. Conversely, let $(U_\alpha)$ be an open covering of
$X$ and, for each $\alpha$, let $u_\alpha:U_\alpha\to Y$ be an $S$-morphism;
if, for every pair of indices $(\alpha,\beta)$, the restrictions of
$u_\alpha$ and of $u_\beta$ to $U_\alpha\cap U_\beta$ coincide, there exists
one and only one $S$-morphism $X\to Y$ whose restriction to each $U_\alpha$
is $u_\alpha$.

If $u:X\to Y$ is an $S$-morphism such that $u(X)\subset V$, where $V$ is an
open subset of $Y$, then $u$, regarded as a morphism of $X$ into $V$, is
again an $S$-morphism.
\end{env}

\begin{env}[2.5.4]
\label{I.2.5.4}
Let $S'\to S$ be a morphism of preschemes; for every $S'$-prescheme, the
composite morphism $X\to S'\to S$ defines $X$ as an $S$-prescheme.
Conversely, suppose that $S'$ is the prescheme induced on an open subset of
$S$; let $X$ be an $S$-prescheme, and suppose that the structural morphism
$f:X\to S$ is such that $f(X)\subset S'$; then $X$ can be regarded as an
$S'$-prescheme. In this latter case, if $Y$ is an $S$-prescheme whose
structural morphism also maps the underlying space into $S'$, every
$S$-morphism of $X$ into $Y$ is also an $S'$-morphism.
\end{env}

\begin{env}[2.5.5]
\label{I.2.5.5}
If $X$ is an $S$-prescheme and $\varphi:X\to S$ the structural morphism, we
call an $S$-\emph{section of $X$} an $S$-morphism of $S$ into $X$, that is
to say, a morphism of preschemes $\psi:S\to X$ such that
$\varphi\circ\psi$ is the identity of $S$. We denote by $\Gamma(X/S)$
the set of $S$-sections of $X$.
\end{env}
